% Options for packages loaded elsewhere
\PassOptionsToPackage{unicode}{hyperref}
\PassOptionsToPackage{hyphens}{url}
\PassOptionsToPackage{dvipsnames,svgnames,x11names}{xcolor}
\documentclass[
  10pt,
  fontset=fandol]{ctexart}
\usepackage{xcolor}
\usepackage[margin=16mm]{geometry}
\usepackage{amsmath,amssymb}
\setcounter{secnumdepth}{-\maxdimen} % remove section numbering
\usepackage{iftex}
\ifPDFTeX
  \usepackage[T1]{fontenc}
  \usepackage[utf8]{inputenc}
  \usepackage{textcomp} % provide euro and other symbols
\else % if luatex or xetex
  \usepackage{unicode-math} % this also loads fontspec
  \defaultfontfeatures{Scale=MatchLowercase}
  \defaultfontfeatures[\rmfamily]{Ligatures=TeX,Scale=1}
\fi
\usepackage{lmodern}
\ifPDFTeX\else
  % xetex/luatex font selection
\fi
% Use upquote if available, for straight quotes in verbatim environments
\IfFileExists{upquote.sty}{\usepackage{upquote}}{}
\IfFileExists{microtype.sty}{% use microtype if available
  \usepackage[]{microtype}
  \UseMicrotypeSet[protrusion]{basicmath} % disable protrusion for tt fonts
}{}
\makeatletter
\@ifundefined{KOMAClassName}{% if non-KOMA class
  \IfFileExists{parskip.sty}{%
    \usepackage{parskip}
  }{% else
    \setlength{\parindent}{0pt}
    \setlength{\parskip}{6pt plus 2pt minus 1pt}}
}{% if KOMA class
  \KOMAoptions{parskip=half}}
\makeatother
\setlength{\emergencystretch}{3em} % prevent overfull lines
\providecommand{\tightlist}{%
  \setlength{\itemsep}{0pt}\setlength{\parskip}{0pt}}
\usepackage{amsmath,amssymb,mathtools,needspace}
\xeCJKDeclareCharClass{CJK}{"2032,"0400->"04FF,"2160->"216B,"2460->"2473,"25A0,"25C7,"25CB,"25CF}
\IfFontExistsTF{Arial Unicode MS}{\xeCJKsetup{AutoFallBack=true}\setCJKfallbackfamilyfont{\CJKrmdefault}{Arial Unicode MS}}{}
\setcounter{secnumdepth}{0}
\setlength{\emergencystretch}{2em}
\allowdisplaybreaks
\usepackage{bookmark}
\IfFileExists{xurl.sty}{\usepackage{xurl}}{} % add URL line breaks if available
\urlstyle{same}
\hypersetup{
  pdftitle={FWT 的设计思想},
  colorlinks=true,
  linkcolor={Maroon},
  filecolor={Maroon},
  citecolor={Blue},
  urlcolor={Blue},
  pdfcreator={LaTeX via pandoc}}

\title{FWT 的设计思想}
\author{}
\date{}

\begin{document}
\maketitle

\subsubsection{10.4.1　FWT 的设计思想}\label{fwt-ux7684ux8bbeux8ba1ux601dux60f3}

在具体设计快速算法 FWT 之前，首先考察两种简单情形。由于 1 阶和 2 阶 Hadamard 阵为

\[
H_1=[\,+\,],
\]

\[
H_2=\begin{pmatrix}+&+\\+&-\end{pmatrix},
\]

因而 1-WT 具有极其简单的形式

\[
X(0)=x(0).
\]

这里输入数据即为所求结果，因而不需要做任何计算。此外，2-WT 为

\[
\begin{cases}
X(0)=x(0)+x(1),\\
X(1)=x(0)-x(1).
\end{cases}
\]

这项计算也很平凡，不存在算法设计问题。

可见，1-WT 与 2-WT 都是极为简单的。

\textbf{快速算法 FWT 的设计思想是，基于规模减半的二分手续，通过 2-WT 的反复计算，将所给 \(N\)-WT 逐步加工成 1-WT，从而得出所求的结果。}

快速算法 FWT 是优秀算法的一朵奇葩，它鲜明地展现了``简单的重复生成复杂''这一算法设计的基本理念。此外，它可以充当一个样板，示范运用二分演化机制设

\href{../../source-images/pdf-271.jpeg}{原页}

计快速变换的全过程。

\begin{center}\rule{0.5\linewidth}{0.5pt}\end{center}

\subsection{相关原页校注（agent 补充）}\label{ux76f8ux5173ux539fux9875ux6821ux6ce8agent-ux8865ux5145}

以下校注来自本节覆盖的原页，可能也涉及同页相邻小节；原文照录与校正说明分开保留。

\subsubsection{原书 PDF 页 270 的校注}\label{ux539fux4e66-pdf-ux9875-270-ux7684ux6821ux6ce8}

\href{../pages/pdf-270.md}{完整原页转录} · \href{../../source-images/pdf-270.jpeg}{原页图像}

校注记录：CH10-E002。

\begin{quote}
校注（agent 补充，CH10-E002）：原文``对称正交阵''照录。依本章未归一化的 \(\pm1\) 元素定义，\(H_N^\mathsf{T}H_N=NI\)，\(N>1\) 时并非通常定义下满足 \(Q^\mathsf{T}Q=I\) 的正交矩阵；准确地说其行列两两正交，\(H_N/\sqrt N\) 才是正交矩阵。原页逆变换中的 \(1/N\) 因子与该定义一致。本项是原书术语疑误，不是 OCR 错误。
\end{quote}

\subsubsection{原书 PDF 页 271 的校注}\label{ux539fux4e66-pdf-ux9875-271-ux7684ux6821ux6ce8}

\href{../pages/pdf-271.md}{完整原页转录} · \href{../../source-images/pdf-271.jpeg}{原页图像}

校注记录：CH10-E001。

\begin{quote}
校注（agent 补充，CH10-E001）：本页化简式及其后两个 \(N/2\)-WT 的第二式，原图均写作 \(H_{N/2}(N/2+i,j)\)，上文已照录。按前页定理 10.1 和本章对 \(N/2\) 阶矩阵的索引范围，这里的行指标 \(N/2+i\) 超出 \(0,\ldots,N/2-1\)；由 \(H_N(N/2+i,j)=H_{N/2}(i,j)\)，两处应为 \(H_{N/2}(i,j)\)。这是原书下标疑误，不是 OCR 错误。\href{../assets/ch10/erratum-p271-indices.png}{原式放大裁片}。此外，本页``定理 1''照录，所用递推关系对应前页``定理 10.1''。
\end{quote}

\subsection{本节覆盖原页的校注记录}\label{ux672cux8282ux8986ux76d6ux539fux9875ux7684ux6821ux6ce8ux8bb0ux5f55}

下列记录按原页关联，含同页相邻小节的校注；计算时同时读取上文校注和完整原页。

\begin{itemize}
\tightlist
\item
  \texttt{CH10-E001}：原书 PDF 页 271
\item
  \texttt{CH10-E002}：原书 PDF 页 270
\end{itemize}

\end{document}
